% GENERATED FILE. Edit canonical lessons, then run npm run generate:books.
% canonical-source-hash: fnv1a64:3fbcc53589fc1b6a
% canonical-lessons: SA-C320-praja, SA-C320-dasah, SA-C320-viplavah, SA-C320-niskasayati, SA-C320-svatantryam

\chapter{Rebellion}
\label{ch:sa-rebellion}
\hlchaptermodality{320}

\begin{chapteropening}
\textbf{By the end of this chapter:} \emph{I can tell how people rose against a ruler.}

Complete the last lesson of chapter 320: I can tell how people rose against a ruler.
\end{chapteropening}

\section[\texorpdfstring{\sk{प्रजा} (prajā)}{प्रजा (prajā)}]{\sk{प्रजा} (prajā) — the subjects of a king; the people}

\label{lesson:SA-C320-praja}



\begin{quote}
Before the new one: say the Sanskrit for retreats, then the Sanskrit for overpowers.
\end{quote}

\subsection*{You'll want to know: \sk{प्रजा}}
\textbf{\sk{प्रजा}} — \emph{prajā} — \textquotedblleft{}the subjects of a king; the people\textquotedblright{}.

\begin{quote}
\sk{प्रजाः} \sk{राज्ये} \sk{सुखेन} \sk{वसन्ति।}
\end{quote}

\begin{quote}
\emph{prajāḥ rājye sukhena vasanti.}
\end{quote}

\begin{quote}
The subjects live happily in the kingdom.
\end{quote}

A feminine noun.

\begin{cousinweb}
\textbf{\sk{प्रजा}} (\emph{prajā}) is built from \emph{pra} and the root √jan, \textquotedblleft{}to be born\textquotedblright{}; MW gives \emph{prajā} as offspring, and also as people, the subjects of a prince (MW).
\end{cousinweb}

\begin{grammarlens}[title={What you've built}]
One more word for telling what happened.
\end{grammarlens}

\subsection*{Your turn}
\begin{itemize}
\raggedright
  \item \emph{Say it:} \emph{prajā}
  \item \emph{Say it:} \emph{prajā}, once more
  \item \emph{Say it:} \emph{prajāḥ rājye sukhena vasanti.} — The subjects live happily in the kingdom.
  \item \emph{Recall:} say the Sanskrit for retreats, then the Sanskrit for overpowers, then say \emph{prajā} again
\end{itemize}

\begin{tcolorbox}[breakable,colback=teal!4,colframe=teal!35!black,title={Before you move on}]
What does \sk{प्रजा} mean? (\textquotedblleft{}The subjects of a king; the people\textquotedblright{}.) Say it once more.
\end{tcolorbox}
\section[\texorpdfstring{\sk{दासः} (dāsaḥ)}{दासः (dāsaḥ)}]{\sk{दासः} (dāsaḥ) — a slave; a servant}

\label{lesson:SA-C320-dasah}



\begin{quote}
Before the new one: say the Sanskrit for overpowers, then the Sanskrit for the subjects.
\end{quote}

\subsection*{You'll want to know: \sk{दासः}}
\textbf{\sk{दासः}} — \emph{dāsaḥ} — \textquotedblleft{}a slave; a servant\textquotedblright{}.

\begin{quote}
\sk{दासः} \sk{क्षेत्रे} \sk{कार्यं} \sk{करोति।}
\end{quote}

\begin{quote}
\emph{dāsaḥ kṣetre kāryaṃ karoti.}
\end{quote}

\begin{quote}
The slave works in the field.
\end{quote}

A masculine noun.

\begin{cousinweb}
\textbf{\sk{दासः}} (\emph{dāsaḥ}) MW's oldest senses are a fiend and the name of certain foes in the Veda, and it also gives \emph{dāsa} as a slave or servant; MW gives no root in the entry, so its origin is left uncertain here (MW).
\end{cousinweb}

\begin{grammarlens}[title={What you've built}]
One more word for telling what happened.
\end{grammarlens}

\subsection*{Your turn}
\begin{itemize}
\raggedright
  \item \emph{Say it:} \emph{dāsaḥ}
  \item \emph{Say it:} \emph{dāsaḥ}, once more
  \item \emph{Say it:} \emph{dāsaḥ kṣetre kāryaṃ karoti.} — The slave works in the field.
  \item \emph{Recall:} say the Sanskrit for overpowers, then the Sanskrit for the subjects, then say \emph{dāsaḥ} again
\end{itemize}

\begin{tcolorbox}[breakable,colback=teal!4,colframe=teal!35!black,title={Before you move on}]
What does \sk{दासः} mean? (\textquotedblleft{}A slave; a servant\textquotedblright{}.) Say it once more.
\end{tcolorbox}
\section[\texorpdfstring{\sk{विप्लवः} (viplavaḥ)}{विप्लवः (viplavaḥ)}]{\sk{विप्लवः} (viplavaḥ) — a revolt; an uprising}

\label{lesson:SA-C320-viplavah}



\begin{quote}
Before the new one: say the Sanskrit for the subjects, then the Sanskrit for a slave.
\end{quote}

\subsection*{You'll want to know: \sk{विप्लवः}}
\textbf{\sk{विप्लवः}} — \emph{viplavaḥ} — \textquotedblleft{}a revolt; an uprising\textquotedblright{}.

\begin{quote}
\sk{राज्ये} \sk{विप्लवः} \sk{अभवत्।}
\end{quote}

\begin{quote}
\emph{rājye viplavaḥ abhavat.}
\end{quote}

\begin{quote}
A revolt broke out in the kingdom.
\end{quote}

A masculine noun.

\begin{cousinweb}
\textbf{\sk{विप्लवः}} (\emph{viplavaḥ}) is built from \emph{vi} and the root √plu, \textquotedblleft{}to float, drift\textquotedblright{}; MW gives \emph{viplava} as confusion, trouble, disaster, and also tumult, affray, revolt (MW).
\end{cousinweb}

\begin{grammarlens}[title={What you've built}]
One more word for telling what happened.
\end{grammarlens}

\subsection*{Your turn}
\begin{itemize}
\raggedright
  \item \emph{Say it:} \emph{viplavaḥ}
  \item \emph{Say it:} \emph{viplavaḥ}, once more
  \item \emph{Say it:} \emph{rājye viplavaḥ abhavat.} — A revolt broke out in the kingdom.
  \item \emph{Recall:} say the Sanskrit for the subjects, then the Sanskrit for a slave, then say \emph{viplavaḥ} again
\end{itemize}

\begin{tcolorbox}[breakable,colback=teal!4,colframe=teal!35!black,title={Before you move on}]
What does \sk{विप्लवः} mean? (\textquotedblleft{}A revolt; an uprising\textquotedblright{}.) Say it once more.
\end{tcolorbox}
\section[\texorpdfstring{\sk{निष्कासयति} (niṣkāsayati)}{निष्कासयति (niṣkāsayati)}]{\sk{निष्कासयति} (niṣkāsayati) — drives out; expels}

\label{lesson:SA-C320-niskasayati}



\begin{quote}
Before the new one: say the Sanskrit for a slave, then the Sanskrit for a revolt.
\end{quote}

\subsection*{You'll want to know: \sk{निष्कासयति}}
\textbf{\sk{निष्कासयति}} — \emph{niṣkāsayati} — \textquotedblleft{}drives out; expels\textquotedblright{}.

\begin{quote}
\sk{प्रजाः} \sk{क्रूरं} \sk{सेनापतिं} \sk{निष्कासयन्ति।}
\end{quote}

\begin{quote}
\emph{prajāḥ krūraṃ senāpatiṃ niṣkāsayanti.}
\end{quote}

\begin{quote}
The people drive out the cruel general.
\end{quote}

Past (he …): \sk{सः} \sk{शत्रुं} \sk{नगरात्} \sk{निष्कासितवान्}, \emph{saḥ śatruṃ nagarāt niṣkāsitavān}.

\begin{cousinweb}
\textbf{\sk{निष्कासयति}} (\emph{niṣkāsayati}) is built from \emph{niṣ}, \textquotedblleft{}out\textquotedblright{}, and the causative of the root √kas, \textquotedblleft{}to go\textquotedblright{}; MW gives \emph{niṣ-kāsayati} as to drive or turn out, expel (MW).
\end{cousinweb}

\begin{grammarlens}[title={What you've built}]
One more word for telling what happened.
\end{grammarlens}

\subsection*{Your turn}
\begin{itemize}
\raggedright
  \item \emph{Say it:} \emph{niṣkāsayati}
  \item \emph{Say it:} \emph{niṣkāsayati}, once more
  \item \emph{Say it:} \emph{prajāḥ krūraṃ senāpatiṃ niṣkāsayanti.} — The people drive out the cruel general.
  \item \emph{Recall:} say the Sanskrit for a slave, then the Sanskrit for a revolt, then say \emph{niṣkāsayati} again
\end{itemize}

\begin{tcolorbox}[breakable,colback=teal!4,colframe=teal!35!black,title={Before you move on}]
What does \sk{निष्कासयति} mean? (\textquotedblleft{}Drives out; expels\textquotedblright{}.) Say it once more.
\end{tcolorbox}
\section[\texorpdfstring{\sk{स्वातन्त्र्यम्} (svātantryam)}{स्वातन्त्र्यम् (svātantryam)}]{\sk{स्वातन्त्र्यम्} (svātantryam) — freedom; independence}

\label{lesson:SA-C320-svatantryam}



\begin{quote}
Before the new one: say the Sanskrit for a revolt, then the Sanskrit for drives out.
\end{quote}

\subsection*{You'll want to know: \sk{स्वातन्त्र्यम्}}
\textbf{\sk{स्वातन्त्र्यम्}} — \emph{svātantryam} — \textquotedblleft{}freedom; independence\textquotedblright{}.

\begin{quote}
\sk{प्रजाः} \sk{स्वातन्त्र्यम्} \sk{इच्छन्ति।}
\end{quote}

\begin{quote}
\emph{prajāḥ svātantryam icchanti.}
\end{quote}

\begin{quote}
The people want freedom.
\end{quote}

A neuter noun.

\begin{cousinweb}
\textbf{\sk{स्वातन्त्र्यम्}} (\emph{svātantryam}) is built on \emph{sva-tantra}, \textquotedblleft{}self-dependent\textquotedblright{} (from \emph{sva}, \textquotedblleft{}own\textquotedblright{}, and \emph{tantra}, \textquotedblleft{}rule\textquotedblright{}), with the lengthened first vowel; MW gives \emph{svātantrya} as following one's own will, independence (MW).
\end{cousinweb}

\begin{grammarlens}[title={What you've built}]
One more word for telling what happened.
\end{grammarlens}

\subsection*{Your turn}
\begin{itemize}
\raggedright
  \item \emph{Say it:} \emph{svātantryam}
  \item \emph{Say it:} \emph{svātantryam}, once more
  \item \emph{Say it:} \emph{prajāḥ svātantryam icchanti.} — The people want freedom.
  \item \emph{Recall:} say the Sanskrit for a revolt, then the Sanskrit for drives out, then say \emph{svātantryam} again
\end{itemize}

\begin{tcolorbox}[breakable,colback=teal!4,colframe=teal!35!black,title={Before you move on}]
What does \sk{स्वातन्त्र्यम्} mean? (\textquotedblleft{}Freedom; independence\textquotedblright{}.) Say it once more.
\end{tcolorbox}
